% Chunk: \S4.4 Structure of the Interaction
% Source: Weinberg QFT Vol. I, physical pp. 205--211 (printed pp. 182--188)
% Status: source-reviewed and compile-clean; equations (4.4.1)--(4.4.10); no figures; no uncertainties

\subsection{Structure of the Interaction}

We now ask, what sort of Hamiltonian will yield an $S$-matrix that satisfies the
cluster decomposition principle? It is here that the formalism of creation and
annihilation operators comes into its own. The answer is contained in the
theorem that the $S$-matrix satisfies the cluster decomposition principle if
(and as far as I know, only if) the Hamiltonian can be expressed as in Eq.
\eqref{eq:4.2.8}:
\begin{align}
H={}&\sum_{N=0}^{\infty}\sum_{M=0}^{\infty}
 \int dq'_1\cdots dq'_N\,dq_1\cdots dq_M \notag\\
&\quad\cdot a^\dagger_{q'_1}\cdots a^\dagger_{q'_N}
 a_{q_M}\cdots a_{q_1} \notag\\
&\quad\cdot h_{NM}(q'_1\cdots q'_N,q_1\cdots q_M).
\tag{4.4.1}
\label{eq:4.4.1}
\end{align}
with coefficient functions $h_{NM}$ that contain just a \emph{single}
three-dimensional momentum-conservation delta function (returning here briefly
to a more explicit notation)
\begin{align}
&h_{NM}(\mathbf{p}'_1\sigma'_1 n'_1\cdots
 \mathbf{p}'_N\sigma'_N n'_N,
 \mathbf{p}_1\sigma_1 n_1\cdots \mathbf{p}_M\sigma_M n_M)
 \notag\\
&\quad={}
 \delta^3(\mathbf{p}'_1+\cdots+\mathbf{p}'_N
 -\mathbf{p}_1-\cdots-\mathbf{p}_N) \notag\\
&\qquad\cdot
 \widetilde h_{NM}(\mathbf{p}'_1\sigma'_1 n'_1\cdots
 \mathbf{p}'_N\sigma'_N n'_N,
 \mathbf{p}_1\sigma_1 n_1\cdots \mathbf{p}_M\sigma_M n_M),
\tag{4.4.2}
\label{eq:4.4.2}
\end{align}
where $\widetilde h_{NM}$ contains no delta function factors. Note that
Eq.~\eqref{eq:4.4.1} by itself has no content---we saw in Section 4.2 that \emph{any}
operator can be put in this form. It is only Eq.~\eqref{eq:4.4.1} combined with the
requirement that $h_{NM}$ has only the single delta function shown in
Eq.~\eqref{eq:4.4.2} that guarantees that the $S$-matrix satisfies the cluster decomposition
principle.

The validity of this theorem in perturbation theory will become obvious when we
develop the Feynman diagram formalism in Chapter 6. The trusting reader may
prefer to skip the rest of the present chapter, and move on to consider the
implications of this theorem in Chapter 5. However, the proof has some
instructive features, and will help to clarify in what sense the field theory of
the next chapter is inevitable.

To prove this theorem, we make use of perturbation theory in its time-dependent
form. (One of the advantages of time-dependent perturbation theory is that it
makes the combinatorics underlying the cluster decomposition principle much
more transparent; if $E$ is a sum of one-particle energies then $e^{-iEt}$ is a
product of functions of the individual energies, while
$[E-E_\alpha+i\epsilon]^{-1}$ is not.) The $S$-matrix is given by Eq.~\eqref{eq:3.5.10}
as\footnote{We are now adopting the convention that for $n=0$, the time-ordered
product in Eq.~\eqref{eq:4.4.3} is taken as the unit operator, so the $n=0$ term in the
sum just yields the term $\delta(\beta-\alpha)$ in $S_{\beta\alpha}$.}
\begin{equation}
S_{\beta\alpha}=\sum_{n=0}^{\infty}\frac{(-i)^n}{n!}
 \int_{-\infty}^{\infty}dt_1\cdots dt_n\,
 \bra{\beta}T\{H_I(t_1)\cdots H_I(t_n)\}\ket{\alpha},
\tag{4.4.3}
\label{eq:4.4.3}
\end{equation}
where the Hamiltonian is split into a free-particle part $H_0$ and an
interaction $H_{\mathrm{int}}$, and
\begin{equation}
H_I(t)\equiv \exp(iH_0t)H_{\mathrm{int}}\exp(-iH_0t).
\tag{4.4.4}
\label{eq:4.4.4}
\end{equation}
Now, the states $\ket{\alpha}$ and $\ket{\beta}$ may be expressed as in
Eq.~\eqref{eq:4.2.2} as products of creation operators acting on the vacuum
$\ket{\Omega}$, and $H_I(t)$ is itself a sum of products of creation and
annihilation operators, so each term in the sum~\eqref{eq:4.4.3} may be written as a sum
of vacuum expectation values of products of creation and annihilation
operators. By using the commutation or anticommutation relations~\eqref{eq:4.2.5} we may
move each annihilation operator in turn to the right past all the creation
operators. For each annihilation operator moved to the right past a creation
operator we have two terms, as shown by writing Eq.~\eqref{eq:4.2.5} in the form
\[
a_{q'}a^\dagger_q=\mathbin{\pm}a^\dagger_q a_{q'}+\delta(q'-q).
\]

Moving other creation operators past the annihilation operator in the first
term generates yet more terms. But Eq.~\eqref{eq:4.2.4} shows that any annihilation
operator that moves all the way to the right and acts on $\ket{\Omega}$ gives
zero, so in the end all we have left is the delta functions. In this way, the
vacuum expectation value of a product of creation and annihilation operators is
given by a sum of different terms, each term equal to a product of delta
functions and $\pm$ signs from the commutators or anticommutators. It follows
that each term in Eq.~\eqref{eq:4.4.3} may be expressed as a sum of terms, each term
equal to a product of delta functions and $\pm$ signs from the commutators or
anticommutators and whatever factors are contributed by $H_I(t)$, integrated
over all the times and integrated and summed over the momenta, spins, and
species in the arguments of the delta functions.

Each of the terms generated in this way may be symbolized by a diagram. (This
is not yet the full Feynman diagram formalism, because we are not yet going to
associate numerical quantities with the ingredients in the diagrams; we are
using the diagrams here only as a way of keeping track of three-momentum delta
functions.) Draw $n$ points, called \emph{vertices}, one for each $H_I(t)$
operator. For each delta function produced when an annihilation operator in one
of these $H_I(t)$ operators moves past a creation operator in the initial state
$\ket{\alpha}$, draw a line coming into the diagram from below that ends at the
corresponding vertex. For each delta function produced when an annihilation
operator in the adjoint of the final state $\ket{\beta}$ moves past a creation
operator in one of the $H_I(t)$, draw a line from the corresponding vertex
upwards out of the diagram. For each delta function produced when an
annihilation operator in one $H_I(t)$ moves past a creation operator in another
$H_I(t)$ draw a line between the two corresponding vertices. Finally, for each
delta function produced when an annihilation operator in the adjoint of the
final state moves past a creation operator in the initial state, draw a line
from bottom to top, right through the diagrams. Each of the delta functions
associated with one of these lines enforces the equality of the momentum
arguments of the pair of creation and annihilation operators represented by the
line. There is also at least one delta function contributed by each of the
vertices, which enforces the conservation of the total three-momentum at the
vertex.

Such a diagram may be connected (every point connected to every other by a set
of lines) and if not connected, it breaks up into a number of connected pieces.
The $H_I(t)$ operator associated with a vertex in one connected component
effectively commutes with the $H_I(t)$ associated with any vertex in any other
connected component, because for this diagram, we are not including any terms
in which an annihilation operator in one vertex destroys a particle that is
produced by a creation operator in the other vertex---if we did, the two
vertices would be in the same connected component. Thus the matrix element in
Eq.~\eqref{eq:4.4.3} can be expressed as a sum over \emph{products} of contributions,
one from each connected component:
\begin{align}
&\bra{\beta}T\{H_I(t_1)\cdots H_I(t_n)\}\ket{\alpha}
\notag\\
&\qquad={}
\sum_{\text{clusterings}}(\pm)\prod_{j=1}^{\nu}
 \Bigl(\bra{\beta_j}T\{H_I(t_{j1})\cdots H_I(t_{jn_j})\}
 \ket{\alpha_j}\Bigr)_C.
\tag{4.4.5}
\label{eq:4.4.5}
\end{align}
Here the sum is over all ways of splitting up the incoming and outgoing
particles and $H_I(t)$ operators into $\nu$ clusters (including a sum over
$\nu$ from 1 to $n$) with the $n_j$ operators
$H_I(t_{j1})\cdots H_I(t_{jn_j})$ and the subsets of initial particles
$\alpha_j$ and final particles $\beta_j$ all in the $j$th cluster. Of course,
this means that
\[
n=n_1+\cdots+n_\nu
\]
and also the set $\alpha$ is the union of all the particles in the subsets
$\alpha_1,\alpha_2,\ldots,\alpha_\nu$, and likewise for the final state. Some
of the clusters in Eq.~\eqref{eq:4.4.5} may contain no vertices at all, i.e., $n_j=0$;
for these factors, we must take the matrix element factor in Eq.~\eqref{eq:4.4.5} to
vanish unless $\beta_j$ and $\alpha_j$ are both one-particle states (in which
case it is just a delta function $\delta(\alpha_j-\beta_j)$), because the only
connected diagrams without vertices consist of a single line running through
the diagram from bottom to top. Most important, the subscript $C$ in
Eq.~\eqref{eq:4.4.5} means that we exclude any contributions corresponding to disconnected
diagrams, that is, any contributions in which any $H_I(t)$ operator or any
initial or final particle is not connected to every other by a sequence of
particle creations and annihilations.

Now let us use Eq.~\eqref{eq:4.4.5} in the sum~\eqref{eq:4.4.3}. Every time variable is integrated
from $-\infty$ to $+\infty$, so it makes no difference which of the
$t_1,\ldots,t_n$ are sorted out into each cluster. The sum over clusterings
therefore yields a factor $n!/n_1!n_2!\cdots n_\nu!$, equal to the number of
ways of sorting out $n$ vertices into $\nu$ clusters, each containing
$n_1,n_2,\ldots$ vertices:
\begin{align*}
&\int_{-\infty}^{\infty}dt_1\cdots dt_n\,
 \bra{\beta}T\{H_I(t_1)\cdots H_I(t_n)\}\ket{\alpha}
\\
&\quad={}
\frac{n!}{n_1!n_2!\cdots n_\nu!}
\sum_{\mathrm{PART}}(\pm)
\sum_{n_1+\cdots+n_\nu=n}
\prod_{j=1}^{\nu}\int_{-\infty}^{\infty}
 dt_{j1}\cdots dt_{jn_j}
\\
&\qquad\quad\cdot
 \Bigl(\bra{\beta_j}T\{H_I(t_{j1})\cdots H_I(t_{jn_j})\}
 \ket{\alpha_j}\Bigr)_C.
\end{align*}
The first sum here is over all ways of partitioning the particles in the initial
and final states into clusters $\alpha_1\cdots\alpha_\nu$ and
$\beta_1\cdots\beta_\nu$ (including a sum over the number $\nu$ of clusters).
The factor $n!$ here cancels the $1/n!$ in Eq.~\eqref{eq:4.4.3}, and the factor $(-i)^n$
in the perturbation series for~\eqref{eq:4.4.5} can be written as a product
$(-i)^{n_1}\cdots(-i)^{n_\nu}$, so instead of summing over $n$ and then summing
separately over $n_1,\ldots,n_\nu$, constrained by
$n_1+\cdots+n_\nu=n$, we can simply sum independently over each
$n_1,\ldots,n_\nu$. This gives finally
\begin{align*}
S_{\beta\alpha}
={}&\sum_{\mathrm{PART}}(\pm)\prod_{j=1}^{\nu}
 \sum_{n_j=0}^{\infty}\frac{(-i)^{n_j}}{n_j!}
 \int_{-\infty}^{\infty}dt_{j1}\cdots dt_{jn_j}
\\
&\quad\cdot
 \Bigl(\bra{\beta_j}T\{H_I(t_{j1})\cdots H_I(t_{jn_j})\}
 \ket{\alpha_j}\Bigr)_C.
\end{align*}
Comparing this with the definition~\eqref{eq:4.2.2} of the connected matrix elements
$S^C_{\beta\alpha}$, we see that these matrix elements are just given by the
factors in the product here
\begin{equation}
S^C_{\beta\alpha}=\sum_{n=0}^{\infty}\frac{(-i)^n}{n!}
 \int_{-\infty}^{\infty}dt_1\cdots dt_n\,
 \Bigl(\bra{\beta}T\{H_I(t_1)\cdots H_I(t_n)\}
 \ket{\alpha}\Bigr)_C.
\tag{4.4.6}
\label{eq:4.4.6}
\end{equation}
(The subscript $j$ is dropped on all the $t$s and $n$s, as these are now mere
integration and summation variables.) We see that $S^C_{\beta\alpha}$ is
calculated by a very simple prescription: \emph{$S^C_{\beta\alpha}$ is the sum
of all contributions to the $S$-matrix that are connected, in the sense that we
drop all terms in which any initial or final particle or any operator $H_I(t)$
is not connected to all the others by a sequence of particle creations and
annihilations.} This justifies the adjective `connected' for $S^C$.

As we have seen, momentum is conserved at each vertex and along every line, so
the connected parts of the $S$-matrix individually conserve momentum:
$S^C_{\beta\alpha}$ contains a factor $\delta^3(\mathbf{p}_\beta-
\mathbf{p}_\alpha)$. What we want to prove is that $S^C_{\beta\alpha}$ contains
no other delta functions.

We now make the assumption that the coefficient fractions $h_{NM}$ in the
expansion~\eqref{eq:4.4.1} of the Hamiltonian in terms of creation and annihilation
operators are proportional to a \emph{single} three-dimensional delta function,
that ensures momenta conservation. This is automatically true for the
free-particle Hamiltonian $H_0$, so it is also then true separately for the
interaction $H_{\mathrm{int}}$. Returning to the graphical interpretation of
the matrix elements that we have been using, this means that each vertex
contributes one three-dimensional delta function. (The other delta functions in
matrix elements $(H_{\mathrm{int}})_{\gamma\delta}$ simply keep the momentum of
any particle that is not created or annihilated at the corresponding vertex
unchanged.) Now, most of these delta functions simply go to fix the momentum of
intermediate particles. The only momenta that are left unfixed by such delta
functions are those that circulate in loops of internal lines. (Any line which
if cut leaves the diagram disconnected carries a momentum that is fixed by
momentum conservation as some linear combination of the momenta of the lines
coming into or going out of the diagram. If the diagram has $L$ lines that can
all be cut at the same time without the diagram becoming disconnected, then we
say it has $L$ independent loops, and there are $L$ momenta that are not fixed
by momentum conservation.) With $V$ vertices, $I$ internal lines, and $L$
loops, there are $V$ delta functions, of which $I-L$ go to fix internal
momenta, leaving $V-I+L$ delta functions relating the momenta of incoming or
outgoing particles. But a well-known topological identity\footnote{A graph
consisting of a single vertex has $V=1$, $L=0$, and $C=1$. If we add $V-1$
vertices with just enough internal lines to keep the graph connected, we have
$I=V-1$, $L=0$, and $C=1$. Any additional internal lines attached (without new
vertices) to the same connected graph produce an equal number of loops, so
$I=V+L-1$ and $C=1$. If a disconnected graph consists of $C$ such connected
parts, the sums of $I$, $V$, and $L$ in each connected part will then satisfy
$\sum I=\sum V+\sum L-C$.} tells that for any graph consisting of $C$
connected pieces, the numbers of vertices, internal lines, and loops are
related by
\begin{equation}
V-I+L=C.
\tag{4.4.7}
\label{eq:4.4.7}
\end{equation}
Hence for a connected matrix element like $S^C_{\beta\alpha}$, which arises
from graphs with $C=1$, we find just a single three-dimensional delta function
$\delta^3(\mathbf{p}_\beta-\mathbf{p}_\alpha)$, as was to be proved.

It was not important in the above argument that the time variables were
integrated from $-\infty$ to $+\infty$. Thus exactly the same arguments can be
used to show that if the coefficients $h_{N,M}$ in the Hamiltonian contain just
single delta functions, then $U(t,t_0)$ can also be decomposed into connected
parts, each containing a single momentum-conservation delta function factor.
On the other hand, the connected part of the $S$-matrix also contains an
energy-conservation delta function, and when we come to Feynman diagrams in
Chapter 6 we shall see that $S^C_{\beta\alpha}$ contains only a single
energy-conservation delta function, $\delta(E_\beta-E_\alpha)$, while
$U(t,t_0)$ contains no energy-conservation delta functions at all.

It should be emphasized that the requirement that $h_{NM}$ in Eq.~\eqref{eq:4.4.1}
should have only a single three-dimensional momentum conservation delta
function factor is very far from trivial, and has far-reaching implications.
For instance, assume that $H_{\mathrm{int}}$ has non-vanishing matrix elements
between two-particle states. Then Eq.~\eqref{eq:4.4.1} must contain a term with $N=M=2$,
and coefficient
\begin{equation}
v_{2,2}(\mathbf{p}'_1\mathbf{p}'_2,
 \mathbf{p}_1\mathbf{p}_2)
=(H_{\mathrm{int}})_{\mathbf{p}'_1\mathbf{p}'_2;
 \mathbf{p}_1\mathbf{p}_2}.
\tag{4.4.8}
\label{eq:4.4.8}
\end{equation}
(We are here temporarily dropping spin and species labels.) But then the matrix
element of the interaction between three-particle states is
\begin{align}
(H_{\mathrm{int}})_{\mathbf{p}'_1\mathbf{p}'_2\mathbf{p}'_3;
 \mathbf{p}_1\mathbf{p}_2\mathbf{p}_3}
={}&v_{3,3}(\mathbf{p}'_1\mathbf{p}'_2\mathbf{p}'_3,
 \mathbf{p}_1\mathbf{p}_2\mathbf{p}_3) \notag\\
&+v_{2,2}(\mathbf{p}'_1\mathbf{p}'_2,
 \mathbf{p}_1\mathbf{p}_2)\,
 \delta^3(\mathbf{p}'_3-\mathbf{p}_3)
 \mathbin{\pm}\text{ permutations}.
\tag{4.4.9}
\label{eq:4.4.9}
\end{align}

As mentioned at the beginning of this chapter, we might try to make a
relativistic quantum theory that is not a field theory by choosing $v_{2,2}$ so
that the two-body $S$-matrix is Lorentz-invariant, and adjusting the rest of the
Hamiltonian so that there is no scattering in states containing three or more
particles. We would then have to take $v_{3,3}$ to cancel the other terms in
Eq.~\eqref{eq:4.4.9}
\begin{align}
v_{3,3}(\mathbf{p}'_1\mathbf{p}'_2\mathbf{p}'_3,
 \mathbf{p}_1\mathbf{p}_2\mathbf{p}_3)
={}&-v_{2,2}(\mathbf{p}'_1\mathbf{p}'_2,
 \mathbf{p}_1\mathbf{p}_2)\,
 \delta^3(\mathbf{p}'_3-\mathbf{p}_3)
 \mathbin{\mp}\text{ permutations}.
\tag{4.4.10}
\label{eq:4.4.10}
\end{align}
However, this would mean that each term in $v_{3,3}$ contains \emph{two} delta
function factors (recall that
$v_{2,2}(\mathbf{p}'_1\mathbf{p}'_2,\mathbf{p}_1\mathbf{p}_2)$ has a factor
$\delta^3(\mathbf{p}'_1+\mathbf{p}'_2-\mathbf{p}_1-\mathbf{p}_2)$) and this
would violate the cluster decomposition principle. Thus in a theory satisfying
the cluster decomposition principle, the existence of scattering processes
involving two particles makes processes involving three or more particles
inevitable.

\begin{center}
$*\ *\ *$
\end{center}

When we set out to solve three-body problems in quantum theories that satisfy
the cluster decomposition principle, the term $v_{3,3}$ in Eq.~\eqref{eq:4.4.9} gives no
particular trouble, but the extra delta function in the other terms makes the
Lippmann--Schwinger equation difficult to solve directly. The problem is that
these delta functions make the kernel
$[E_\alpha-E_\beta+i\epsilon]^{-1}(H_{\mathrm{int}})_{\beta\alpha}$ of this
equation not square-integrable, even after we factor out an overall momentum
conservation delta function. In consequence, it cannot be approximated by a
finite matrix, even one of very large rank. To solve problems involving three
or more particles, it is necessary to replace the Lippmann--Schwinger equation
with one that has a connected right-hand side. Such equations have been
developed for the scattering of three or more particles~\hyperref[ref:4.8]{[8,}\hyperref[ref:4.9]{9]}, and in
non-relativistic scattering problems they can be solved recursively, but they
have not turned out to be useful in relativistic theories and so will not be
described in detail here.

However, recasting the Lippmann--Schwinger equation in this manner is useful in
another way. Our arguments in this section have so far relied on perturbation
theory. I do not know of any non-perturbative proof of the main theorem of this
section, but it has been shown~\hyperref[ref:4.9]{[9]} that these reformulated non-perturbative
dynamical equations are \emph{consistent} with the requirement that
$U^C(t,t_0)$ (and hence $S^C$) should also contain only a single
momentum-conservation delta function, as required by the cluster decomposition
principle, provided that the Hamiltonian satisfies our condition that the
coefficient functions $h_{N,M}$ each contain only a single
momentum-conservation delta function.
